\documentclass[12pt, reqno]{amsart}
\usepackage{latexsym, amssymb, amscd, xypic, amsthm, setspace}
\input xy
\xyoption{all}

%%: Font = times
\usepackage{mathptmx} % rm & math
\usepackage[scaled=0.90]{helvet} % ss
\usepackage{courier} % tt
\normalfont
\usepackage[T1]{fontenc}
\pagestyle{plain}


\newtheorem{theorem}{Theorem}[subsection]
\newtheorem{prop}[theorem]{Proposition}
\newtheorem{lemma}[theorem]{Lemma}
\newtheorem{corollary}[theorem]{Corollary}
\newtheorem{claim}[theorem]{Claim}

\theoremstyle{definition}
\newtheorem{de}[theorem]{Definition}
\newtheorem{notation}{Notation}
\newtheorem{example}[theorem]{Example}
\newtheorem{examples}{Examples}
\newtheorem{remark}[theorem]{Remark}
\newtheorem{remarks}{Remarks}
\newtheorem{question}{Question}


\newtheorem*{nthm}{Theorem}
\newtheorem*{nlem}{Lemma}
\newtheorem*{nprop}{Proposition}
\newtheorem*{ncor}{Corollary}
\newtheorem*{nclaim}{Claim}
\newtheorem*{nde}{Definition}
\newtheorem*{nex}{Example}

\newtheoremstyle{test}{3pt}{3pt}{\itshape}{\parindent}{\bfseries}{---}{  }{\thmnote{#3}} \theoremstyle{test}
\newtheorem*{mythm}{}


%: paragraphs

\newcommand{\point}{\vspace{3mm}\par \noindent (\refstepcounter{subsection}{\bf \thesubsection) } }

\newcommand{\tpoint}[1]{\vspace{3mm}\par \noindent \refstepcounter{subsubsection}{\bf \thesubsubsection.}
  {\em #1. ---} }
\newcommand{\epoint}[1]{\vspace{3mm}\par \noindent \refstepcounter{subsection}{\bf \thesubsection.}
  {\em #1.} }
\newcommand{\spoint}{\vspace{3mm}\par \noindent \refstepcounter{subsection}{\bf \thesubsection.} }

\newcommand{\bpoint}[1]{\vspace{3mm}\par \noindent \refstepcounter{subsection}{\bf \thesubsection.}
  {\bf #1.} }



%: numbering and margins

\numberwithin{equation}{subsection}

\setlength{\oddsidemargin}{0cm} \setlength{\evensidemargin}{0cm}
\setlength{\marginparwidth}{0in}
\setlength{\marginparsep}{0in}
\setlength{\marginparpush}{0in}
\setlength{\topmargin}{0in}
\setlength{\headheight}{0pt}
\setlength{\headsep}{0pt}
\setlength{\footskip}{.3in}
\setlength{\textheight}{9.2in}
\setlength{\textwidth}{6.5in}
\setlength{\parskip}{0.25pt}
\setlength{\parindent}{0.25in}

%: math shortcuts


\DeclareMathOperator{\inv}{inv}
\DeclareMathOperator{\out}{Out}
\DeclareMathOperator{\gal}{Gal}
\renewcommand{\mod}{\mathrm{\, mod \, }}
\renewcommand{\d}[1]{{}^{#1}}
\renewcommand{\b}[1]{\mathbf{{#1}}}
\newcommand{\C}{\mathbb{C}}
\newcommand{\be}[1]{\begin{eqnarray} \label{#1}}
\newcommand{\ee}{\end{eqnarray}}
\newcommand{\re}{\mathsf{Re}}
\newcommand{\im}{\mathsf{Im}}
\newcommand{\rr}{\rightarrow}

\newcommand{\n}{\mathbb{N}}
\newcommand{\zee}{\mathbb{Z}}
\newcommand{\R}{\mathbb{R}}
\newcommand{\pr}{\mathbb{P}}
\renewcommand{\gcd}[2]{\mathsf{gcd}(#1, #2)}
\newcommand{\ov}[1]{\overline{#1}}
\renewcommand{\Re}{\mathrm{Re}}
\renewcommand{\Im}{\mathrm{Im}}



\usepackage{enumitem}

\title{Midterm 1 Review Sheet \\ MATH 411: Honors Complex Variables, Fall 2026}
\date{last compiled: \today}

\begin{document}
\maketitle

\begin{center}\fbox{\begin{minipage}{0.95\textwidth}\vspace{4pt}\centering
\begin{tabular}{rl}
\textbf{Date:} & Friday, October 9, 2026, in class (CAB 369, 3:00--3:50 PM)\\[2pt]
\textbf{Length:} & 50 minutes\\[2pt]
\textbf{Format:} & Closed book. No notes, no textbook, no calculators or devices.\\[2pt]
\textbf{Coverage:} & Lectures 1--12; Lang \S\S1.1--1.7, 2.1--2.5, 3.2--3.3 (through Goursat)
\end{tabular}\vspace{4pt}\end{minipage}}\end{center}

\noindent Since the exam is closed book, you should be able to \emph{state} the main
definitions and theorems below precisely, not just recognize them. When you use
a result from class in a solution, state clearly which result you are using and
check that its hypotheses hold.

\noindent Lecture 13 (local primitives, Cauchy's theorem on a disc) and Lecture 14 (computing real
integrals) will \emph{not} be tested on this midterm.

\section*{Topics}

\subsection*{1. Basic complex numbers \normalfont (Lectures 1--3; Lang \S\S1.1--1.4)}
\begin{itemize}[leftmargin=*]
\item Arithmetic in $\C$: real and imaginary parts, the conjugate $\ov z$, the modulus $|z|=\sqrt{z\ov z}$,
      and the identities $|zw|=|z||w|$, $\ov{zw}=\ov z\,\ov w$, $z^{-1}=\ov z/|z|^2$.
\item The triangle inequality $|z+w|\le|z|+|w|$ and its reverse form $\bigl||z|-|w|\bigr|\le|z-w|$.
\item Polar form $z=re^{i\theta}$; multiplication multiplies moduli and adds arguments.
      The definition $e^{a+ib}=e^a(\cos b+i\sin b)$.
\item Solving $z^n=w$: the $n$ distinct $n$-th roots of a nonzero $w$, and roots of unity.
\item Stereographic projection and the extended plane $\widehat\C=\C\cup\{\infty\}$.
\item Topology of $\C$: open and closed discs, open and closed sets, interior, boundary and adherent points;
      limits and continuity of functions $f:S\to\C$ in $\varepsilon$--$\delta$ form.
\end{itemize}

\subsection*{2. $\C$-differentiability \normalfont (Lectures 3--5; Lang \S\S1.5--1.7)}
\begin{itemize}[leftmargin=*]
\item Definition of $f'(z_0)$ as a limit over $h\in\C$, and of \emph{holomorphic} on an open set.
      The equivalent ``linear approximation'' form $f(z_0+h)-f(z_0)=f'(z_0)h+h\,\sigma(h)$ with $\sigma(h)\to0$.
\item Differentiable implies continuous; sum, product, quotient and chain rules.
      Polynomials are entire; rational functions are holomorphic away from the zeros of the denominator.
\item Total (real) differentiability of $F=(u,v):U\subseteq\R^2\to\R^2$; the Jacobian;
      continuous partials $\Rightarrow$ totally differentiable.
\item \textbf{Cauchy--Riemann criterion:} $f=u+iv$ is $\C$-differentiable at $z=x+iy$ if and only if
      $F=(u,v)$ is totally differentiable at $(x,y)$ and $u_x=v_y$, $u_y=-v_x$ there.
      In that case $f'=u_x+iv_x$. Know the proof, not just the statement.
\item Using CR to show a function is (or is not) holomorphic; e.g.\ $e^z$ is entire, $\ov z$ and $|z|^2$ are not holomorphic.
\item On a connected open set, $f'\equiv0$ implies $f$ is constant (via ``locally constant'').
      Typical consequences: if $f$ is holomorphic on a disc and $\Re f$, $\Im f$ or $|f|$ is constant, then $f$ is constant.
\item Conformality: if $f'(z_0)\ne0$, then $f$ preserves oriented angles between curves at $z_0$.
\item Green's theorem and the resulting first version of Cauchy's theorem
      (assuming $f$ has continuous partials).
\end{itemize}

\subsection*{3. Power series \normalfont (Lectures 6--9; Lang \S\S2.1--2.2)}
\begin{itemize}[leftmargin=*]
\item Formal power series $\C[[T]]$: addition, multiplication (Cauchy product),
      inversion ($f$ is invertible iff $a_0\ne0$), composition $f\circ g$ when $g$ has no constant term.
      Be able to compute the first few coefficients of $1/f$, $f\cdot g$, $f\circ g$.
\item Convergence of sequences and series in $\C$; Cauchy sequences; absolute convergence; the comparison test.
\item Sequences and series of functions: pointwise vs.\ uniform convergence, the sup norm $\|\cdot\|_S$,
      uniformly Cauchy sequences. A uniform limit of continuous functions is continuous.
      The comparison test for series of functions (Weierstrass $M$-test).
\item \textbf{Radius of convergence:} every $\sum a_nz^n$ has $0\le r\le\infty$ such that it converges absolutely
      on $D(0,r)$, uniformly on each closed disc $\ov D(0,s)$ with $s<r$, and diverges for $|z|>r$.
\item $\limsup$ and \textbf{Hadamard's formula} $1/r=\limsup_n|a_n|^{1/n}$. Be able to compute radii of convergence
      (including series with gaps, such as $\sum a_nz^{n^2}$), and know the proof.
\item Behaviour on the boundary circle can vary: compare $\sum nz^n$, $\sum z^n/n$, $\sum z^n/n^2$.
\item The formal derivative $\sum na_nz^{n-1}$ has the same radius of convergence, and a power series can be
      differentiated term by term inside its disc of convergence; hence it is infinitely $\C$-differentiable there.
\end{itemize}

\subsection*{4. Analytic functions \normalfont (Lectures 9--11; Lang \S\S2.3--2.5)}
\begin{itemize}[leftmargin=*]
\item Definition: $f$ is \emph{analytic} on $U$ if near each $z_0\in U$ it is given by a convergent power series in $z-z_0$.
\item A power series with radius $r$ is analytic on all of $D(0,r)$ (re-centring a power series).
\item Analytic $\Rightarrow$ holomorphic, in fact infinitely $\C$-differentiable; the Taylor coefficients are
      $a_k=f^{(k)}(z_0)/k!$. Sums and products of analytic functions are analytic.
\item Uniqueness of power series expansions; \textbf{isolated zeros}; the \textbf{identity principle}
      (two convergent power series that agree at a sequence of nonzero points tending to $0$ are the same power series).
\item Analytic extensions of real functions, and uniqueness of such extensions; the complex exponential, $\sin$, $\cos$.
\end{itemize}

\subsection*{5. Path integrals and the Cauchy--Goursat theorem \normalfont (Lectures 11--12; Lang \S\S3.2--3.3)}
\begin{itemize}[leftmargin=*]
\item Integrals of complex-valued functions of a real variable, and $\bigl|\int_a^bF\bigr|\le\int_a^b|F|$.
\item Piecewise $C^1$ paths and the path integral $\int_\gamma f\,dz=\int_a^b f(\gamma(t))\gamma'(t)\,dt$.
      Be able to compute integrals over line segments, circles and rectangles
      (e.g.\ $\int_{|z|=1}z^n\,dz$ for all $n\in\zee$).
\item Independence of parametrization; reversing orientation changes the sign.
\item \textbf{The ML-estimate:} $\bigl|\int_\gamma f\bigr|\le\bigl(\sup_{\gamma}|f|\bigr)\cdot L(\gamma)$.
\item For a polynomial $Q$ and a closed path $\gamma$, $\int_\gamma Q=0$.
\item \textbf{Cauchy--Goursat:} if $f$ is holomorphic on an open set $U$ and $R\subseteq U$ is a closed rectangle, then
      $\int_{\partial R}f=0$. Know the structure of the proof: subdivide into four rectangles,
      pick the worst one, get nested rectangles shrinking to a point $z_0$, use the linear approximation of $f$ at $z_0$
      together with the lemma on polynomials, and finish with the ML-estimate.
      Note that no continuity of $f'$ is assumed (compare with the Green's theorem version).
\end{itemize}

\section*{How to prepare}
\begin{itemize}[leftmargin=*]
\item Work through Problem Sets 1 and 2.
\item Try the 2021 and 2022 Midterm 1 papers (posted with solutions on the course site) under exam conditions: 50 minutes,
      nothing on your desk. Note that the 2022 exam allowed a page of notes; this one does not.
\item Practise writing out the key statements from memory: the definition of $\C$-differentiability,
      the Cauchy--Riemann criterion, Hadamard's formula, the definition of analytic, the identity principle, the ML-estimate,
      and Cauchy--Goursat.
\end{itemize}

\section*{Practice problems}
\noindent These are for practice only; solutions will not be collected.

\begin{enumerate}[leftmargin=*]

\item Find all solutions of $z^4=-16$ and plot them. Then describe geometrically the set
      $\{z\in\C:\ |z-1|=|z+i|\}$.

\item Show that if $|w|<1$ and $|z|<1$, then $\left|\dfrac{z-w}{1-\ov wz}\right|<1$.
      \emph{(Hint: compare $|1-\ov wz|^2-|z-w|^2$.)}

\item Using the Cauchy--Riemann criterion, determine the set of points at which each function is $\C$-differentiable:
\begin{enumerate}
\item $f(z)=|z|^2$;
\item $f(x+iy)=x^2+iy^2$;
\item $f(z)=z\,\ov z^{\,2}$.
\end{enumerate}

\item Let $f$ be holomorphic on the unit disc and suppose that $|f|$ is constant. Prove that $f$ is constant.

\item Find the radius of convergence of
\begin{enumerate}
\item $\displaystyle\sum_{n\ge0}\frac{n^3}{3^n}\,z^n$;
\item $\displaystyle\sum_{n\ge0}2^n z^{2n}$;
\item $\displaystyle\sum_{n\ge1}n!\,z^{n^2}$.
\end{enumerate}

\item Find the terms of order $\le3$ in the power series of $\dfrac{e^z}{1-z}$ at $z=0$, and of
      $\dfrac{1}{\cos z}$ at $z=0$. What is the radius of convergence of the first?

\item Show directly from the definition that $f(z)=1/z$ is analytic on $\C\setminus\{0\}$:
      expand it as a power series around an arbitrary $z_0\ne0$ and find the radius of convergence.

\item Let $f$ be analytic on the unit disc $D(0,1)$.
\begin{enumerate}
\item If $f(1/n)=1/n^2$ for all $n\ge2$, determine $f$.
\item Is there such an $f$ with $f(1/n)=(-1)^n/n$ for all $n\ge2$? Justify your answer.
\end{enumerate}

\item Suppose $f(z)=\sum a_nz^n$ has radius of convergence $r>0$ and $f(z)=f(-z)$ for all $|z|<r$. Show that $a_n=0$ for all odd $n$.

\item Compute $\displaystyle\int_\gamma \ov z\,dz$ where $\gamma$ is the unit circle traversed once counterclockwise, and also
      where $\gamma$ is the boundary of the square with vertices $0,1,1+i,i$. Why does this not contradict the Cauchy--Goursat theorem?

\item Let $C_R$ denote the circle $|z|=R$ traversed once counterclockwise, with $R>1$. Use the ML-estimate to show
\[
\left|\int_{C_R}\frac{dz}{z^2+1}\right|\le\frac{2\pi R}{R^2-1},
\]
and conclude that the integral tends to $0$ as $R\to\infty$.

\item State the Cauchy--Goursat theorem carefully. In its proof, explain (i) why the nested rectangles shrink to a single point
      $z_0\in R$, (ii) exactly where the hypothesis that $f$ is $\C$-differentiable at $z_0$ is used, and (iii) how the
      estimate $\bigl|\int_{\partial R}f\bigr|\le 4^n\bigl|\int_{\partial R_n}f\bigr|$ combines with the ML-estimate to finish the argument.

\end{enumerate}

\end{document}
